\documentclass[border=6pt]{standalone}
\usepackage{amsmath}
\usepackage{amssymb}
\usepackage{array}
\usepackage{booktabs}
\begin{document}
\begin{tabular}{@{} >{\raggedright\arraybackslash}p{5cm} >{\centering\arraybackslash}p{7cm} @{}}
\toprule
\multicolumn{2}{c}{\textbf{CONCEPT SUMMARY: Solving Rational Equations}} \\
\midrule
\textbf{Words} & \textbf{Algebra} \\
\midrule
A rational equation is an equation that contains a rational expression. To solve, identify the domain for the variable. Then multiply both sides of the equation by a common denominator and solve. An extraneous solution is a solution that is not valid because that value is excluded from the domain of the original equation.
&
\(\begin{gathered}
\dfrac{1}{x}+\dfrac{2}{x}=\dfrac{1}{6} \quad \text{Domain: } x\ne 0 \\[6pt]
6x\left(\dfrac{1}{x}+\dfrac{2}{x}\right)=6x\left(\dfrac{1}{6}\right) \\
6+12=x \\
18=x \\[4pt]
\text{The domain includes } x=18\text{, so the solution to the equation is } 18.
\end{gathered}\)
\\[10pt]
\cline{2-2}
&
\(\begin{gathered}
\dfrac{x^{2}+4}{x-1}=\dfrac{5}{x-1} \quad \text{Domain: } x\ne 1 \\[6pt]
(x-1)\left(\dfrac{x^{2}+4}{x-1}\right)=(x-1)\left(\dfrac{5}{x-1}\right) \\
x^{2}+4=5 \\
x^{2}=1 \\
x=\pm 1 \\[4pt]
\text{The domain does not include } x=1\text{, so } 1 \text{ is an extraneous solution.} \\
\text{It does include } x=-1\text{, so the solution to the equation is } -1.
\end{gathered}\)
\\
\bottomrule
\end{tabular}
\end{document}
